\section{Every experiment on record}

Table~\ref{tab:experiment-index} is parsed from the research record's own
generated index, which is rebuilt from the experiment manifests by
\pathname{make index}. Nothing in it is retyped.

\begin{table}[htbp]
\caption{The nineteen experiments this report draws on. A superseded entry is
kept deliberately, so that a dead end is walked into once rather than twice.
Produced by \pathname{tables/make_tables.py}.}
\label{tab:experiment-index}
\centering\scriptsize
\begin{tabular}{@{}L{0.15\textwidth}L{0.37\textwidth}L{0.13\textwidth}L{0.25\textwidth}@{}}
\toprule
ID & Title & Status & Primary metric \\
\midrule
EXP-2026-07-001 & Arm A --- fully deterministic pipeline (planning, extraction, synthesis) & complete & correctness 0.9972 (n=1) \\
EXP-2026-07-002 & Arm B --- LLM planning with deterministic extraction and an LLM writer & complete & correctness 0.8505 (n=1) \\
EXP-2026-07-003 & Arm C --- authored plan with an LLM writer reading raw evidence & complete & correctness 0.9690 (n=1) \\
EXP-2026-07-004 & Arm D --- authored plan with an LLM writer reading verified claims & complete & correctness 0.9917 (n=1) \\
EXP-2026-07-005 & Arm E --- LLM planning with the facet-chasing retrieval loop enabled & complete & correctness 0.8560 (n=1) \\
EXP-2026-07-006 & Arm F --- authored plan, facet-chasing loop, LLM writer over verified claims & complete & correctness 0.9957 (n=1) \\
EXP-2026-07-007 & Arm G --- generic primitives only, every hand-written template disabled & complete & correctness 0.6388 (n=1) \\
EXP-2026-07-008 & Component-level evaluation of the planner and both query specialists & complete & planner\_\allowbreak source\_\allowbreak recall 1.0000 (n=1) \\
EXP-2026-07-009 & End-to-end evaluation against the four goals.md pipeline metrics & complete & cross\_\allowbreak source\_\allowbreak answer\_\allowbreak accuracy 1.0000 (n=1) \\
EXP-2026-07-011 & Portability audit --- what survives a dataset that is not MITRE & complete & --- \\
EXP-2026-08-001 & Is Arm G's gap a capability gap in the schema-driven layer, or domain knowledge? & complete & gold\_\allowbreak in\_\allowbreak evidence 1.0000 (n=1) \\
EXP-2026-08-002 & Does the deterministic pipeline's near-ceiling correctness survive a blind, uniformly-sampled holdout? & complete & correctness (not recorded) \\
EXP-2026-08-003 & Do the schema-driven primitives generalise to questions nobody wrote templates for? & complete & gold\_\allowbreak in\_\allowbreak evidence 0.8952 (n=1) \\
EXP-2026-08-004 & Does making the cross-source bridge an outer join close the residual blind-set gap without regressing? & complete & blind\_\allowbreak gold\_\allowbreak in\_\allowbreak evidence 1.0000 (n=1) \\
EXP-2026-08-005 & Is the blind holdout's directional loss a property of these templates, or of a missing schema-driven check? & complete & blind\_\allowbreak gold\_\allowbreak in\_\allowbreak evidence 0.8952 (n=1) \\
EXP-2026-08-006 & Does the outer-join bridge close the residual blind gap in the default cascade without regressing it? & complete & blind\_\allowbreak gold\_\allowbreak in\_\allowbreak evidence 1.0000 (n=1) \\
EXP-2026-08-007 & Judged, at last --- does the repaired cascade answer better, and do the generic and LLM-loop arms differ from it on anything that matters? & complete & blind\_\allowbreak correctness\_\allowbreak cascade 0.8367 (n=30) \\
EXP-2026-08-008 & Do the answer-writing fixes help, how do four arms compare on unseen questions, and does an LLM writing its own queries beat templates? & complete & llm\_\allowbreak query\_\allowbreak vs\_\allowbreak templates\_\allowbreak fresh\_\allowbreak bank (not recorded) \\
EXP-2026-07-010 & First independent cross-model grading pass (the pre-fix committed baseline) & superseded $\rightarrow$ EXP-2026-07-001 & correctness 0.8946 (n=1) \\
\bottomrule
\end{tabular}

\end{table}

\section{The claims ledger as it stands today}

Table~\ref{tab:claims} is parsed from \pathname{docs/research/CLAIMS.md}. Claim
text is truncated; the ledger carries the full wording and the evidence trail.

\begin{table}[htbp]
\caption{Claim status at 2026-08-20, straight from the ledger. Produced by
\pathname{tables/make_tables.py}.}
\label{tab:claims}
\centering\scriptsize
\begin{tabular}{@{}L{0.08\textwidth}L{0.62\textwidth}L{0.18\textwidth}@{}}
\toprule
ID & Claim & Status \\
\midrule
C-001 & On an anchored, in-distribution question set, a fully deterministic pipeline --- template routing, deterministic extraction, deterministic\,\ldots & partial \\
C-002 & Substituting LLM query planning for an authored plan is the dominant single source of correctness loss in the pipeline, with extraction and synthesis\,\ldots & partial \\
C-003 & The hand-written query templates, not the planner, are where the system's accuracy lives: disabling them costs roughly three times what substituting\,\ldots & supported \\
C-004 & The deterministic path is not merely competitive with the LLM-planned path but orders of magnitude cheaper: the full 120-question catalogue completes\,\ldots & supported \\
C-005 & Binding an LLM narrator to the \emph{verified claim set} rather than to raw retrieved evidence removes the accuracy-versus-readability trade-off. & supported \\
C-006 & The value of an iterative retrieval loop is a function of plan quality: it compensates for weak planning at substantial cost and does not substitute\,\ldots & supported \\
C-007 & Constructing queries by filling validated read-only templates places query validity and safety adherence at a ceiling by construction, rather than\,\ldots & supported \\
C-008 & A deterministic safety gate over templated queries blocks adversarial probes without leakage. & partial \\
C-009 & Self-judged evaluation materially overstates answer quality relative to independent cross-model judging, which is why the primary evaluation must be\,\ldots & supported \\
C-010 & The system's dataset-specific surface is bounded and countable rather than pervasive: the schema-agnostic machinery transfers, and re-authoring cost\,\ldots & supported \\
C-012 & Routing to a fixed library of validated parameterised templates outperforms letting an LLM plan and generate the queries, in correctness and in cost. & supported \\
C-017 & Claim-level verification with enforced abstention carrying executed-query proof eliminates unsupported claims without reducing answer coverage. & partial \\
C-021 & Evaluating a graph-question-answering system on questions co-authored with the code that answers them overstates its correctness by a measurable and\,\ldots & supported \\
C-011 & Federated question answering across cybersecurity knowledge graphs held in \emph{separate engines with different query languages} is a contribution\,\ldots & unsupported \\
C-013 & The near-ceiling correctness of the deterministic path is a genuine capability rather than an artifact of a question set co-authored with the\,\ldots & refuted \\
C-014 & The coverage cliff of template routing is measurable: the design's \emph{price} on open-ended, non-anchored questions is bounded and reportable, not\,\ldots & partial \\
C-015 & Federating across two engines outperforms merging both sources into a single graph and querying that. & unsupported \\
C-016 & Mechanically localising the last pipeline stage at which the gold answer was still recoverable contradicts the assumption that LLM pipeline failures\,\ldots & partial \\
\bottomrule
\end{tabular}

\end{table}

Two rows in that table disagree with their own notes and should be reconciled
before anything is written from it.

C-003 reads \texttt{supported} while its notes record that the no-templates arm
judged level with the cascade and that the claim needs rewriting. On the
evidence in \S4.4 the status is at best \texttt{partial}.

C-012 reads \texttt{supported} on correctness. Its wording also claims a cost
advantage, and no cost figure was taken in the run that supports it.

\section{Reproducing the numbers in this report}

One script extracts every value in the figures and tables from the experiments
into one file. The figure and table generators read only that file:

\begin{quote}\small\ttfamily
python3 data/collect.py \\
python3 figures/make\_figures.py \\
python3 tables/make\_tables.py
\end{quote}

\pathname{data/collect.py} raises rather than defaulting when a key is absent,
so a figure cannot be built from a metric that has gone missing. It reads each
experiment's \pathname{data/metrics.json} and the second night's
\pathname{docs/evaluation/night2-2026-08-20/metrics.json}.

Raw run artifacts are not in the research record. They live under
\pathname{docs/evaluation/}, and each experiment carries a directory hash in its
\pathname{data/POINTERS.md} that folds every file's path and content digest in
sorted order.

\section{Known problems with this report}

\begin{description}
  \item[The second night is not yet an experiment.] Section 4.6, 4.7, the second
    generated question set and the token-budget re-run come from a sweep output
    directory rather than from a scaffolded experiment. By this project's own
    rules those numbers do not yet count. This report quotes them anyway,
    labelled, on the grounds that a briefing that omitted them would be less
    honest than one that flags them.
  \item[Two graders are mixed across the record.] Table~\ref{tab:seven-arm} and
    the first blind measurement come from a Sonnet~5 panel. Everything from
    2026-08-19 onward comes from a locally served grader calibrated against it.
    Figure~\ref{fig:six-numbers} uses the local grader throughout for that
    reason, which is why its first blind bar reads 0.7200 rather than the
    Sonnet panel's 0.689.
  \item[Status.] \textbf{Draft.} No second person has worked through the
    checklist. Reviewer and date: none.
\end{description}
